\documentclass[11pt,a4paper]{article}

\usepackage[utf8]{inputenc}
\usepackage[T1]{fontenc}

\IfFileExists{spanish.ldf}{%
  \usepackage[spanish,es-noshorthands,es-tabla]{babel}%
}{%
  \usepackage[english]{babel}%
  \AtBeginDocument{\def\today{\number\day\ de \ifcase\month\or enero\or
    febrero\or marzo\or abril\or mayo\or junio\or julio\or agosto\or
    septiembre\or octubre\or noviembre\or diciembre\fi\ de \number\year}}%
}

\IfFileExists{XCharter.sty}{\usepackage{XCharter}}{\usepackage{charter}}
\IfFileExists{inconsolata.sty}{\usepackage[scaled=0.96]{inconsolata}}{\usepackage{courier}}

\usepackage[a4paper,left=2.4cm,right=2.4cm,top=2.4cm,bottom=2.5cm,headsep=14pt]{geometry}
\usepackage{microtype}
\usepackage{xcolor}
\usepackage{booktabs}
\usepackage{array}
\usepackage{multirow}
\usepackage{float}
\usepackage{enumitem}
\usepackage{caption}
\usepackage{titlesec}
\usepackage{fancyhdr}
\usepackage{tocloft}
\usepackage{listings}
\usepackage{parskip}
\usepackage[hidelinks]{hyperref}

\definecolor{gris}{gray}{0.42}
\definecolor{grisclaro}{gray}{0.72}

\linespread{1.09}
\raggedbottom
\renewcommand{\arraystretch}{1.3}

\AtBeginDocument{%
  \renewcommand{\contentsname}{Contenido}%
  \renewcommand{\tablename}{Tabla}%
  \renewcommand{\lstlistingname}{Código}%
}

\captionsetup{font=small,labelfont=bf,labelsep=period,skip=8pt}

\titleformat{\section}
  {\normalfont\large\bfseries}{\llap{\thesection\hspace{0.7em}}}{0pt}{}
\titlespacing*{\section}{0pt}{2.3\baselineskip}{0.7\baselineskip}

\titleformat{\subsection}
  {\normalfont\normalsize\bfseries}{\llap{\thesubsection\hspace{0.7em}}}{0pt}{}
\titlespacing*{\subsection}{0pt}{1.5\baselineskip}{0.35\baselineskip}

\renewcommand{\cfttoctitlefont}{\large\bfseries}
\renewcommand{\cftsecfont}{\bfseries}
\renewcommand{\cftsecpagefont}{\bfseries}
\renewcommand{\cftsubsecpagefont}{\color{gris}}
\renewcommand{\cftsecleader}{\hfill}
\renewcommand{\cftsubsecleader}{\color{grisclaro}\cftdotfill{\cftdotsep}}
\setlength{\cftbeforesecskip}{9pt}
\setlength{\cftaftertoctitleskip}{14pt}
\setlength{\cftsubsecindent}{1.6em}

\pagestyle{fancy}
\fancyhf{}
\renewcommand{\headrulewidth}{0.4pt}
\renewcommand{\headrule}{\hbox to\headwidth{\color{grisclaro}\leaders\hrule height \headrulewidth\hfill}}
\fancyhead[L]{\footnotesize\color{gris}\textls[40]{TRABAJO FINAL}}
\fancyhead[R]{\footnotesize\color{gris}\thepage}
\makeatletter
\let\ps@plain\ps@fancy
\makeatother

\lstset{
  basicstyle=\ttfamily\small,
  keywordstyle=\bfseries,
  commentstyle=\itshape\color{gris},
  stringstyle=\itshape,
  numbers=left,
  numberstyle=\ttfamily\scriptsize\color{gris},
  numbersep=12pt,
  stepnumber=1,
  frame=single,
  framerule=0.4pt,
  rulecolor=\color{grisclaro},
  framesep=8pt,
  xleftmargin=26pt,
  framexleftmargin=18pt,
  xrightmargin=8pt,
  breaklines=true,
  breakatwhitespace=true,
  showstringspaces=false,
  tabsize=2,
  columns=fullflexible,
  keepspaces=true,
  captionpos=b,
  aboveskip=1.4em,
  belowskip=1.4em,
  extendedchars=true,
  literate=%
    {á}{{\'a}}1 {é}{{\'e}}1 {í}{{\'i}}1 {ó}{{\'o}}1 {ú}{{\'u}}1
    {Á}{{\'A}}1 {É}{{\'E}}1 {Í}{{\'I}}1 {Ó}{{\'O}}1 {Ú}{{\'U}}1
    {ñ}{{\~n}}1 {Ñ}{{\~N}}1 {ü}{{\"u}}1 {Ü}{{\"U}}1
}

\newcommand{\pordefinir}{\textit{\color{gris}[por definir]}}
\newcommand{\colhead}[1]{\footnotesize\bfseries\textls[50]{\MakeUppercase{#1}}}
\newcommand{\sep}{\hspace{0.5em}\textcolor{grisclaro}{$\cdot$}\hspace{0.5em}}

\begin{document}

\begin{center}
  {\huge\bfseries Trabajo Final}

  \vspace{7pt}

  {\large\color{gris}Bases de Datos 1}

  \vspace{16pt}

  Jerónimo Hoyos Botero\sep Dylan Ramírez Blanquicet
\end{center}

\vspace{1.1cm}

\tableofcontents
\thispagestyle{empty}

\newpage

\section{Esquema escogido para la implementación}
\label{sec:esquema}

De las relaciones obtenidas en el trabajo anterior se seleccionaron cuatro entidades
conectadas entre sí: \textbf{Especie}, \textbf{Pokémon}, \textbf{Entrenador Pokémon}
y \textbf{Tipo de Objeto}.

\pordefinir\ Descripción de los atributos, cardinalidades, opcionalidades y decisiones
de diseño tomadas para las cuatro relaciones.

\section{Fase I: Implementación de esquemas}
\label{sec:fase1}

Los esquemas se implementaron en MySQL y PostgreSQL con todas las restricciones
definidas: clave primaria, claves foráneas, \textit{checks}, no nulos y unicidad.

\subsection{Creación de tablas en MySQL}

\begin{lstlisting}[language=SQL,caption={Definición del esquema en MySQL.},label={lst:mysql}]
CREATE TABLE especie (
);
\end{lstlisting}

\subsection{Creación de tablas en PostgreSQL}

\begin{lstlisting}[language=SQL,caption={Definición del esquema en PostgreSQL.},label={lst:postgres}]
CREATE TABLE especie (
);
\end{lstlisting}

\section{Fase II: Carga de datos}
\label{sec:fase2}

Se realizaron tres pruebas de inserción en cada SGBD, con 1.000, 10.000 y 100.000
filas por tabla, midiendo el tiempo empleado en cada caso.

\subsection{Código de población}

\begin{lstlisting}[language=SQL,caption={Población de las bases de datos.},label={lst:poblacion}]

\end{lstlisting}

\subsection{Tiempos de inserción}

\begin{table}[H]
  \centering
  \caption{Tiempo de inserción por prueba de carga, en segundos.}
  \label{tab:insercion}
  \begin{tabular}{@{}lcc@{}}
    \toprule
    \colhead{Prueba de carga} & \colhead{MySQL} & \colhead{PostgreSQL} \\
    \midrule
    \addlinespace[2pt]
    Carga 1 \textcolor{gris}{(1.000 filas)}   & \pordefinir & \pordefinir \\
    Carga 2 \textcolor{gris}{(10.000 filas)}  & \pordefinir & \pordefinir \\
    Carga 3 \textcolor{gris}{(100.000 filas)} & \pordefinir & \pordefinir \\
    \addlinespace[2pt]
    \bottomrule
  \end{tabular}
\end{table}

\section{Fase III: Consultas en SQL}
\label{sec:fase3}

\subsection{Consulta 1: JOIN entre tres tablas}

\pordefinir\ Descripción del objetivo de la consulta.

\begin{lstlisting}[language=SQL,caption={Consulta 1.},label={lst:consulta1}]
SELECT
FROM
;
\end{lstlisting}

\subsection{Consulta 2: GROUP BY sobre dos tablas, JOIN con otra y ORDER BY}

\pordefinir\ Descripción del objetivo de la consulta.

\begin{lstlisting}[language=SQL,caption={Consulta 2.},label={lst:consulta2}]
SELECT
FROM
GROUP BY
ORDER BY
;
\end{lstlisting}

\section{Fase IV: Ejecución de las consultas en Python}
\label{sec:fase4}

Desde Python se establece conexión con cada SGBD, se traen las tablas completas
mediante \texttt{SELECT * FROM tabla} hacia estructuras de datos y se resuelven las
dos consultas con código propio, midiendo el tiempo de solución.

\subsection{Conexión y medición de tiempos}

\begin{lstlisting}[language=Python,caption={Conectividad y medición de tiempos.},label={lst:conexion}]

\end{lstlisting}

\subsection{Solución de las consultas en Python}

\begin{lstlisting}[language=Python,caption={Resolución de las consultas en Python.},label={lst:pyconsultas}]

\end{lstlisting}

\section{Resultados comparativos}
\label{sec:resultados}

\begin{table}[H]
  \centering
  \caption{Tiempos de respuesta de las consultas por escenario y carga, en segundos.}
  \label{tab:consultas}
  \begin{tabular}{@{}llccc@{}}
    \toprule
    \colhead{Consulta} & \colhead{Carga} & \colhead{MySQL} & \colhead{PostgreSQL} & \colhead{Python} \\
    \midrule
    \addlinespace[2pt]
    \multirow{3}{*}{Consulta 1}
      & 1.000   & \pordefinir & \pordefinir & \pordefinir \\
      & 10.000  & \pordefinir & \pordefinir & \pordefinir \\
      & 100.000 & \pordefinir & \pordefinir & \pordefinir \\
    \addlinespace[4pt]
    \multirow{3}{*}{Consulta 2}
      & 1.000   & \pordefinir & \pordefinir & \pordefinir \\
      & 10.000  & \pordefinir & \pordefinir & \pordefinir \\
      & 100.000 & \pordefinir & \pordefinir & \pordefinir \\
    \addlinespace[2pt]
    \bottomrule
  \end{tabular}
\end{table}

\section{Conclusiones}
\label{sec:conclusiones}

\begin{itemize}[leftmargin=1.3em,itemsep=0.45em,label=\textcolor{gris}{\textbullet}]
  \item \pordefinir\ Análisis de los tiempos de inserción entre MySQL y PostgreSQL.
  \item \pordefinir\ Comparación de los tiempos de consulta entre los SGBDs y Python.
  \item \pordefinir\ Observaciones finales y aprendizajes del trabajo.
\end{itemize}

\end{document}